
\vspace{2mm}
\begin{table*}[ht!]
\centering
\small
\begin{tabular}{p{4.2cm} p{5.8cm} p{5.8cm}}
\toprule
\textbf{Aspect} & \textbf{Strength of GRACE} & \textbf{Limitation / Caution} \\
\midrule
\textbf{Representation Geometry} & Enforces structured clusters for safe/unsafe/jailbreak responses & May require behavior labels or clustering heuristics \\
\textbf{Pooling Strategy} & Learnable attention over LLM layers reveals alignment-relevant depth & Static and prompt-invariant; dynamic variants may help \\
\textbf{Parameter Efficiency} & Only attention weights trained; backbone frozen & May underutilize full model capacity in latent alignment \\
\textbf{Adversarial Robustness} & Reduces ASR by 35–39\%, outperforming DPO by 6–8$\times$ & Assumes adversarial samples are correctly labeled and separable \\
\textbf{Scalability} & Works with any frozen LLM checkpoint & Forward-pass cost increases with number of behavior classes \\
\textbf{Generalization} & Effective across jailbreak, control, and degradation attacks & Not tested on multimodal or instruction-following benchmarks \\
\bottomrule
\end{tabular}
\caption{At-a-glance summary of GRACE’s strengths and limitations.}
\label{tab:grace_summary}
\vspace{-1mm}
\end{table*}

\section{Discussion and Limitations}

\paragraph{Representation-Grounded Alignment.}
GRACE introduces a paradigm shift from output-based preference tuning to geometry-aware alignment, showing that internal representations encode critical safety-relevant information. Our latent contrastive losses reshape the internal geometry of LLMs to reflect structured behavioral distinctions, enforcing compactness within unsafe regions and separation from safe clusters. This alignment of latent geometry boosts adversarial robustness and paves the way for explainable and interpretable safety enforcement.

\paragraph{Latent Contrastive Supervision vs. Traditional Preference Learning.}
While DPO and its variants align model behavior through pairwise preference loss, they overlook the internal mechanisms that lead to unsafe completions. GRACE complements preference learning by supervising these mechanisms directly in the embedding space. Our contrastive losses target adversarial proximity and unsafe dispersion—factors often missed by output-only training. This hybrid formulation leads to sharper representation boundaries and better generalization of unseen attacks.

\paragraph{Efficiency and Interpretability.}
GRACE is highly parameter-efficient: the only trainable parameters during pooling are the scalar layerwise weights $\alpha^{(l)}$. The rest of the model remains frozen during this step, enabling fast convergence and modular analysis. This structure enables post-hoc auditing of layer contributions to alignment and offers an interpretable bridge between model depth and safety fidelity. Furthermore, the pooled representations offer new debugging and safety attribution tools, which can benefit practitioners seeking deeper control over LLM behavior.

\paragraph{Limitations.}
Despite strong empirical results, GRACE has certain limitations:
\begin{itemize}
    \item \textbf{Behavioral triplet assumption:} GRACE operates under a semi-synthetic triplet construction where (safe, unsafe, jailbreak) completions are drawn from separate datasets. This assumption may introduce distributional shifts or confounding signals when true behavior-specific clusters are not well-separated.
    \item \textbf{Frozen backbone constraint:} During contrastive supervision, the LLM is frozen. While this improves modularity and efficiency, it limits the system’s ability to jointly co-adapt latent and output layers for optimal alignment.
    \item \textbf{Static pooling:} The learned attention profile over layers is static and prompt-invariant. Dynamic, prompt-aware or multi-head pooling might further improve semantic disentanglement in future versions.
    \item \textbf{Compute overhead:} Each batch requires multiple forward passes (one per behavior class), marginally increasing compute costs during latent supervision.
    \item \textbf{Modality and dataset limitations:} We evaluate GRACE only on text-based LLMs. Its extension to multimodal models and richer alignment benchmarks (e.g., Anthropic's HH-RLHF or red-teaming datasets) remains an open direction.
\end{itemize}

\paragraph{Future Extensions.}
We envision several promising extensions to GRACE:
\begin{itemize}
    \item \textit{Prompt-conditional attention pooling} for adaptive safety supervision.
    \item \textit{Joint training of latent and policy layers}, allowing end-to-end preference tuning under geometric constraints.
    \item \textit{Geometric alignment diagnostics}, where AVQI and cluster shape are tracked during training to assess overfitting, drift, or compression.
    \item \textit{Multi-agent adversarial alignment}, where GRACE-inspired contrastive losses are used across interacting LLM agents in competitive tasks.
\end{itemize}



\vspace{1mm}
\noindent Overall, GRACE provides a blueprint for bridging latent-space structure and alignment-aware tuning. It invites a broader shift from black-box preference optimization to interpretable, mechanistically grounded fine-tuning of language models.
